\chapter*{TP s'entrainer à utiliser des listes}

\begin{Exercise}[title={Moyenne}]
\Question Écrire une fonction  \textit{moyenne(liste)} qui retourne la moyenne des éléments de  \textit{liste}.
\Question Faites des tests en appelant cette fonction dans le programme principale
\end{Exercise}

\begin{Exercise}[title={Médiane}]
\Question Écrire une fonction  \textit{mediane(liste)} qui retourne la médiane des éléments de  \textit{liste}.
\Question Faites des tests en appelant cette fonction dans le programme principale
\end{Exercise}

\begin{Exercise}[title={Diviseurs d'un entier}]
\Question Écrire une fonction  \textit{diviseurs(n)} qui renvoie la liste des diviseurs positifs de  \textit{n}, dans l'ordre croissant.
\Question Faites des tests en appelant cette fonction dans le programme principale
\end{Exercise}

\begin{Exercise}[title={Étendue}]
\Question Écrire une fonction \textit{etendue(liste)} qui retourne l'étendue (différence entre le max. et le min.) de  \textit{liste}.
\Question Faites des tests en appelant cette fonction dans le programme principale
\end{Exercise}

\begin{Exercise}[title={Doubles}]
\Question Écrire une fonction  \textit{doubler(liste)} qui renvoie la liste des doubles des nombres de  \textit{liste}. 
\Question Faites des tests en appelant cette fonction dans le programme principale
\end{Exercise}

\begin{Exercise}[title={ Somme des éléments d'une liste}]
\Question Écrire une fonction  \textit{somme(liste)} qui retourne la somme des éléments de  \textit{liste}.
\Question Faites des tests en appelant cette fonction dans le programme principale
\end{Exercise}

\begin{Exercise}[title={Produit des éléments d'une liste}]
\Question Écrire une fonction  \textit{produit(liste)} qui retourne le produit des éléments de  \textit{liste}.
\Question Faites des tests en appelant cette fonction dans le programme principale
\end{Exercise}

\begin{Exercise}[title={ Nombre de pairs}]
\Question Écrire une fonction  \textit{nombrePairs(liste)} qui retourne, pour une liste de nombres entiers, le nombre d'entiers pairs de la  \textit{liste}.
\Question Faites des tests en appelant cette fonction dans le programme principale
\end{Exercise}

\begin{Exercise}[title={Nombre de négatifs}]
\Question Écrire une fonction  \textit{nombreNégatifs(liste)} qui retourne le nombre de nombres négatifs de la  \textit{liste}.
\Question Faites des tests en appelant cette fonction dans le programme principale
\end{Exercise}

\begin{Exercise}[title={Transformation}]
\Question Écrire une fonction  \textit{zeroSiNegatif(liste)} qui retourne  \textit{liste} (une liste d'entiers relatifs) modifiée de sorte que les éléments négatifs sont remplacés par 0. 
\Question Faites des tests en appelant cette fonction dans le programme principale
\end{Exercise}